\section{Native implementation and certificate boundaries}
\label{sec:implementation}

\subsection{Core cancellation integration}

The implementation lives in \path{Topology/UnknotRecognition/fast}.
The successful row-leaf fast path is retained. For a rectangular core,
\code{core_handlebody_certificate} now calls
\code{column_cancellation_plan}, which scans immutable source incidence
and returns a schedule together with its residual rows and structural
counters. The old \code{capped_cancellation_plan} remains available for
compatibility and paired comparisons; its completeness documentation is
corrected as in Proposition~\ref{prop:old-complete}.

The new module \code{column_cancellation_verify} supplies two checks:
\begin{itemize}
\item \code{verify_column_cancellations} checks deletion timestamps.
\item \code{verify_raw_cancellation_obstruction} checks the stopping
condition directly.
\end{itemize}
Neither imports a producer.
The core verifier supports a nested
\code{normal-core-handlebody-v2} certificate with method
\code{column_leaf}. Existing \code{normal-core-handlebody-v1} leaf and
capped certificates remain accepted. The surrounding whole-cut certificate
schemas need not change simply because their nested core proof is faster.

The old square row-leaf obstruction path remains a specialized early exit.
The canonical stopping-set theorem applies directly to the public column
planner's returned residual; a legacy row-leaf residual is not silently
relabelled as that canonical source set. Callers requiring a new independent
stopping certificate should use the column planner and the stopping checker.

Linear cancellation complexity is a statement about the kernel after
source incidence has been built. Existing forest validation, source geometry
reconstruction and component assembly keep their own costs. We do not
claim that modifying the planner makes those unrelated operations linear.

\subsection{Source compilation for saturated bases}

The new \code{normal_cut_saturated} module adds a separate whole-cut query.
Its low-level finite operation is \code{saturated_source_census}; the public
wrapper is \code{normal_cut_saturated}. These are opt-in APIs. They do not
dispatch a new verdict in the main knot recognizer.

The finite source compiler performs the following operations.
\begin{enumerate}
\item Choose a forest in the actual core zero/one-handle graph and recover
the core components. A core three-handle triggers an inconclusive result.
\item Read the authenticated base profiles and ordered attachment components.
Reject closed bases. Each ported base must have multiplicity one, and each
maximal attachment must have a unique base owner and a unique core owner.
\item Count complete attaching circles for each base. A base is saturated
exactly when this count equals its total number of boundary components.
It then has no proper interval attachments. Unported boundary circles are
included in the profile count, so they cannot be accidentally ignored.
\item Extract the old two-handle rows from the actual cyclic source flags.
For each saturated base, traverse all of its full annular words and sum
their old nonforest one-handle crossings into one row.
\item Apply complete column cancellation to the combined source matrix.
If it succeeds, compute connected attachment components and their Euler
characteristics, then use the proved handlebody class to obtain genera.
\end{enumerate}

The fourth step needs more than a word of core vertex labels. An individual
core one-handle may be a loop, so its two ends can belong to the same core
zero-handle. The compiler checks the tetrahedron/face flags at both ends
of every face-sector port, retains the direction of the corresponding
midpoint path, and checks that consecutive path ends join to a closed
curve. Every crossing is counted in absolute value, including repeated
passages and passages with opposite orientations. The independent verifier
reconstructs the same geometry with a separate path routine.

\begin{proposition}[Size of the compiled cancellation kernel]
\label{prop:compiled-size}
Let the old core have $f$ two-handles and $c$ nonforest one-handles, and
let $s$ bases be saturated. The compiled matrix has $f+s$ rows and $c$
possible columns. If $I_0$ is the number of nonzero old core incidences
and $P$ is the total number of active ports, it has at most $I_0+P$
nonzeros. In particular, its size is independent of the numerical sum
$\sum_i(1-\chi(F_i))$ of added band handles.
\end{proposition}

\begin{proof}
There is one row per old two-handle and exactly one per saturated base.
Only old nonforest columns are stored. Each full annular boundary word
uses source ports; the attachment skeleton partitions the ports among
maximal attachments. A core one-handle crossing uses such a port, and
summing crossings into rows can merge nonzeros but cannot create more
than the number of contributing ports. Theorem~\ref{thm:saturated-even-columns}
shows that no omitted band column is needed for the cancellation decision.
\end{proof}

In the repository's source representation the finite core and active
ports are bounded in terms of $T$\cite{repo}. Even using the coarse
$I_0\le fc$, this is a polynomial-size kernel. With explicit source
incidence lists, construction is charged to their actual total length.
Binary Euler characteristics and multiplicities are retained as integers
and use the established hexadecimal codec for transport.

\subsection{The whole-cut proof and its independent replay}

A successful wrapper returns a \code{normal-cut-saturated-v1} certificate
with three components: the parallelity source proof, the ordered-boundary
source proof, and a \code{saturated-source-census-v1} finite proof. The
finite proof names the forest, every saturated base and its relation label,
the full cancellation schedule, and the resulting census.

The independent \code{normal_cut_saturated_verify} module first replays
both existing source proofs. Only then does it inspect the source geometry,
base profiles and boundary skeleton contained in them. It independently
checks the forest, reconstructs saturated ownership and raw crossing
counts, checks every pivot, and computes attachment components by graph
traversal. The producer uses a different component-construction path.
Both check the redundant source Euler identity
\[
 \chi(\text{cut})=\chi(\text{core})+\chi(\text{bundles})
                  -\chi(\text{interfaces}).
\]
That identity is an arithmetic consistency check after the geometric proof;
it is not the proof that a component is a handlebody.

The supplied finite saturated verifier deliberately retains a direct scan
of surviving source rows at each pivot. Its matrix-check cost is polynomial,
with a coarse $O(mI+m^2)$ bound, rather than the linear timestamp bound
already implemented in the separate core checker. This keeps the new
geometric replay easy to inspect. Theorem~\ref{thm:timestamp} identifies
a later compatible optimization if this finite replay becomes material.
Whole-cut cost still includes the source orbit proofs, which dominate
many of the measured small examples.

There are six source orbit traces across the parallelity and boundary
queries. The producer's cycle allowance is shared across the two stages;
the wrapper subtracts the consumed parallelity cycles before the boundary
query. Replay counts all six supplied traces against its operation cap.
Caller abort exceptions propagate. Certificate mutation checks reject
invalid schema fields, repeated pivots, nonexistent or Boolean labels,
changed relative-handle counts, altered source charts and changed genus
claims. The existing exact integer codec handles large signed and
hexadecimal quantities.

\subsection{Examples the extension must distinguish}

\paragraph{A saturated disk.}
Here $\chi(F)=1$, so the relative bundle has no one-handles and one
two-handle. Its row is the full annular attachment row. The new construction
recovers the existing disk-base absorption case.

\paragraph{A saturated punctured surface with a core anchor.}
Let a once-bordered orientable genus-$q$ surface be attached to a solid
torus along a core annulus meeting its cocore once. The aggregate matrix
is $(1)$, for every $q$. The old core handle cancels the residual two-handle,
leaving a genus-$2q$ handlebody. The matrix size is constant even if $q$
has thousands of bits. A nonorientable base of crosscap number $q$ gives
the same calculation with a genus-$q$ result.

\paragraph{Two individually primitive annuli on the same handle.}
Take an annulus base and attach both boundary annuli to longitudinal
annuli of one solid torus. The aggregate row is $(2)$, so the criterion
refuses. Treating the two boundary circles as separate singleton rows
would be erroneous. With product gluings, this construction is the product
of a circle with a disk to which one orientable band has been attached;
the result is $T^2\times I$, not a handlebody. Its fundamental group is
$\mathbb Z^2$. The aggregate row detects the missing relative primitivity.

\paragraph{A saturated M\"obius base attached with multiplicity three.}
The vertical annulus of the orientable interval bundle over a M\"obius
band winds twice around its solid-torus core. Gluing it to a winding-three
annulus in another solid torus gives the amalgam
$\langle x,y\mid x^3=y^2\rangle$, the familiar torus-knot group used in
the incoming counterexample\cite{hierarchy-incoming}. The old aggregate
row $(3)$ is blocked. Its infinite cyclic abelianization does not justify
a handlebody conclusion. With a winding-one attachment instead, the row
is $(1)$ and the core generator cancels.

\paragraph{Interleaved core and bundle relations.}
The combined matrix need not be solved in two phases. A saturated row can
expose the next old core row, and an old core cancellation can expose a
saturated row. The producer handles this mixed schedule directly.
If the combined matrix fully cancels, its original-core row submatrix also
fully cancels: any stopping set of that submatrix would obstruct the
combined matrix. Thus joint cancellation cannot cure a fixed old-core
raw obstruction.

The distinction concerns the choice and retention of a compatible
schedule. For the algebraic matrix
\[
 A=\begin{pmatrix}1&1&1\\2&1&0\end{pmatrix},
\]
declare the first row to be an old core relation. The joint column
schedule pivots first in its third column, then in the second row's
second column. But the core-only matrix also admits a pivot in its
first column. Committing to that choice transports the second row to
$(3,2)$, which is blocked. An arbitrary core-only success and its unmarked
census therefore do not replace a schedule checked against all annular
data. This example asserts an algebraic scheduling fact, not an
authenticated source realization of that matrix.

\subsection{Exactly what a successful result certifies}

The output is a complete binary histogram of handlebody genera for a
\emph{supplied normal cut}. Isolated bundle components contribute their full
multiplicities. The result distinguishes one ball from several balls and
from a ball accompanied by positive-genus components.

A knot conclusion needs the separate facts that the input is the authenticated
exterior of the queried diagram and that the cutting data supplies the
appropriate disk or hierarchy argument. The wrapper itself proves neither
of those facts. In particular, a raw stopping certificate proves failure
of a restricted cancellation model, while a resource-limited or unsupported
query returns \code{INCONCLUSIVE}. These are part of the mathematical
input/output contract, not optional interpretations of the results.
